% Macros shared by the printed and the web version of the blueprint.
% Both blueprint/src/web.tex and blueprint/src/print.tex \input this file.
%
% The notation follows Peter Gacs, "Descriptional Complexity and Randomness"
% (arXiv:2105.04704).  The definitions below are reproduced, with the author's
% permission, from his style file gacs.sty and from the preamble of his
% manuscript, restricted to what Chapter 1 actually uses.  Three kinds of
% adjustment were needed, and are marked in place:
%   (1) definitions made with \def are restated with \newcommand,
%   (2) font-switching text macros are reduced to \text,
%   (3) his list environments used package options that plasTeX, which builds
%       the web version, does not implement.
%
% Do not define \lean, \leanok, \uses, \proves, \notready, \discussion or
% \mathlibok here: the web version gets them from the blueprint plasTeX
% plugin and the printed version from macros/print.tex.

% ---------------------------------------------------------------------------
% Theorem environments
% ---------------------------------------------------------------------------
% Gacs runs two independent counters, both reset at every section.  Theorems
% have one of their own; lemmas, propositions, corollaries, definitions,
% remarks and examples share the other.  This is why his text contains both a
% Theorem 1.4.2 and a Corollary 1.4.2.  Reproducing it keeps every reference
% in the blueprint equal to the number printed in his manuscript.
%
% This assumes a document class with chapters (report or book) and a single
% \chapter for Chapter 1, so that the section counter reads 1.3, 1.4 and so on.

\newtheorem{theorem}{Theorem}[section]
\newtheorem{lemma}{Lemma}[section]
\newtheorem{proposition}[lemma]{Proposition}
\newtheorem{corollary}[lemma]{Corollary}

\theoremstyle{definition}
\newtheorem{definition}[lemma]{Definition}
\newtheorem{example}[lemma]{Example}

\theoremstyle{remark}
\newtheorem{remark}[lemma]{Remark}

% Used to flag every place where the Lean statement departs from the printed
% one.  Unnumbered, so it never enters the dependency graph.
\newtheorem*{formalization}{Formalization note}

\theoremstyle{plain}

% ---------------------------------------------------------------------------
% Lists
% ---------------------------------------------------------------------------
% (3) His originals take enumitem key-value options, which plasTeX does not
% implement.  The enumerate package, which it does implement, gives the same
% result.  Caveat: the printed version shows (a), (b) and (i) as intended,
% while plasTeX drops the label option and the web version falls back to
% 1., 2.  Fixing that needs a CSS rule, not a change here.

\usepackage{enumerate}
\newenvironment{alphenum}{\begin{enumerate}[(a)]}{\end{enumerate}}
\newenvironment{romanenum}{\begin{enumerate}[(i)]}{\end{enumerate}}
\newenvironment{bullets}{\begin{itemize}}{\end{itemize}}

% ---------------------------------------------------------------------------
% Number systems and script letters
% ---------------------------------------------------------------------------
\newcommand{\dB}{\mathbb{B}}
\newcommand{\dC}{\mathbb{C}}
\newcommand{\dN}{\mathbb{N}}
\newcommand{\dQ}{\mathbb{Q}}
\newcommand{\dR}{\mathbb{R}}
\newcommand{\dS}{\mathbb{S}}
\newcommand{\dZ}{\mathbb{Z}}

\newcommand{\cF}{\mathcal{F}}
\newcommand{\cG}{\mathcal{G}}
\newcommand{\cI}{\mathcal{I}}
\newcommand{\cM}{\mathcal{M}}
\newcommand{\cQ}{\mathcal{Q}}
\newcommand{\cT}{\mathcal{T}}
\newcommand{\cV}{\mathcal{V}}

% ---------------------------------------------------------------------------
% Greek abbreviations
% ---------------------------------------------------------------------------
\newcommand{\ag}{\alpha}
\newcommand{\bg}{\beta}
\newcommand{\dg}{\delta}
\newcommand{\eps}{\varepsilon}
\newcommand{\ig}{\iota}
\newcommand{\og}{\omega}
\newcommand{\Lg}{\Lambda}
\newcommand{\Og}{\Omega}

% ---------------------------------------------------------------------------
% Relations up to an additive or a multiplicative constant
% ---------------------------------------------------------------------------
% These four are the heart of his notation: a plus or a star under the
% relation means "up to an additive constant" or "up to a constant factor".
\newcommand{\lea}{\stackrel{{}_+}{<}}
\newcommand{\gea}{\stackrel{{}_+}{>}}
\newcommand{\eqa}{\stackrel{{}_+}{=}}
\newcommand{\lem}{\stackrel{{}_*}{<}}
\newcommand{\gem}{\stackrel{{}_*}{>}}
\newcommand{\eqm}{\stackrel{{}_*}{=}}

\newcommand{\sbsq}{\subseteq}
\newcommand{\prefix}{\sqsubseteq}

% ---------------------------------------------------------------------------
% Delimiters and spacing
% ---------------------------------------------------------------------------
\newcommand{\verythinmathskip}{\mskip 1mu}
\newcommand{\setof}[1]{\mathopen\{\verythinmathskip#1\verythinmathskip\mathclose\}}
\newcommand{\ang}[1]{{\mathopen\langle #1\mathclose\rangle}}
\newcommand{\tup}[1]{(#1)}
\newcommand{\pair}[2]{\mathopen(#1,#2\mathclose)}
\newcommand{\tupcod}[1]{\mathopen\langle #1\mathclose\rangle}
\newcommand{\tupdec}[1]{\mathopen[#1\mathclose]}
\newcommand{\cei}[1]{{\lceil #1\rceil}}
\newcommand{\flo}[1]{{\lfloor #1\rfloor}}
\newcommand{\ol}[1]{\overline{#1}}
\newcommand{\card}[1]{|#1|}
\newcommand{\clint}[2]{[#1,#2]}
\newcommand{\opint}[2]{(#1,#2)}
\newcommand{\rint}[2]{(#1,#2]}

% The conditioning bar, as in K(x \mvert y).
\newcommand{\mvert}{\,\vert\,}

% ---------------------------------------------------------------------------
% Complexity-specific notation
% ---------------------------------------------------------------------------
\newcommand{\len}[1]{l(#1)}          % length of a string
\newcommand{\m}{\mathbf{m}}          % (1) his \def\m{\mathbf{m}}
\newcommand{\frto}[2]{#1:#2}         % the range in \xi_{\frto{1}{n}}
\newcommand{\txt}[1]{\text{#1}}      % (2) font switches dropped
\newcommand{\cnv}{\txt{cnv}}
\newcommand{\df}[1]{\emph{#1}}       % his emphasis for a term being defined
\newcommand{\nn}{\nonumber}
